% -----------------------------------------------------------------------------------
\section{Introduction}\label{sec:intro}

Is there a decision theory that gets every decision problem right? For half a century the question
has been pursued case by case. A theory is proposed; a problem is constructed on which its verdict
seems wrong; the theory is amended or a rival proposed; and the rival meets its own
counterexample. The impression this history leaves is that no theory can be optimal everywhere,
because the standards by which theories are judged pull in different directions. This paper tries to
turn that impression into theorems. It asks, for precisely defined criteria of optimality and
precisely defined classes of problems, when a decision theory satisfying them exists and when it
does not, and it tries to close the question in the sense of determining every inconsistency among
the criteria considered and exhibiting a witness for every consistent combination.

\subsection{From dominance to Newcomb's problem}

The expected-utility theories of \textcite{ramsey1931}, \textcite{vonneumannmorgenstern1944}, and
\textcite{savage1954} model an agent's options as functions from states of the world to
consequences, with the states not depending on what the agent does. Under that assumption, the
principle of dominance---prefer an option that is better in every state---is a consequence of
maximising expected utility, and the two principles never conflict. \textcite{jeffrey1965} replaced
the act-independent states with a single algebra of propositions, valuing an act by the news it
would bring; \textcite{bolker1966,bolker1967} supplied the representation theorem. In Jeffrey's theory
the probability of a state may depend on the act, and dominance and expected utility can come apart.

\textcite{nozick1969} made the possibility vivid. A reliable predictor has placed a large sum in an
opaque box if and only if it predicted that the agent would take only that box; a transparent box
contains a small sum. Taking both boxes dominates, since the contents are fixed; taking only the
opaque box maximises expected utility conditional on the act. Nozick credited the problem to the
physicist William Newcomb and proposed restricting dominance to cases where the states are
probabilistically independent of the acts. \textcite{gibbardharper1978} gave the conflict its
canonical form by distinguishing two expected utilities: one computed with the probabilities of
counterfactual conditionals, which measures what an act would causally bring about, and Jeffrey's,
which measures what it indicates. Causal decision theory, as the first came to be called, takes both
boxes; evidential decision theory takes one. Gibbard and Harper also introduced Death in Damascus, in
which Death, having predicted the agent's destination, waits there: whichever city the agent
chooses, the other looks better. \textcite{lewis1979} observed that the prisoner's dilemma played
against a sufficiently similar partner is a Newcomb problem, and \textcite{lewis1981cdt} unified the
causal theories of Gibbard and Harper, \textcite{skyrms1980,skyrms1982}, and Sobel \autocite[see][]{sobel1994} by means of
dependency hypotheses. Against the charge that causal decision theorists end up poorer, Lewis
replied that Newcomb's problem rewards irrationality \autocite{lewis1981rich}. On the evidential side,
the tickle defence of \textcite{eells1982} argued that a self-aware agent's own beliefs and desires screen
off the correlations on which one-boxing relies, and \textcite{spohn1977} argued that deliberation
crowds out prediction, so that an agent cannot coherently assign probabilities to its own current
options. \textcite{joyce1999} gave causal decision theory a representation theorem, and by the end of
the century it had become the orthodox view \autocite{weirich2024sep}.

\subsection{Instability, ratifiability, and the counterexamples}

Death in Damascus exposed a second fault line. \textcite{jeffrey1981,jeffrey1983} proposed that a
rational choice should be \emph{ratifiable}: optimal relative to the beliefs the agent would have
upon making it. In Death in Damascus no pure option is ratifiable. \textcite{harper1986} proposed
maximising causal expected utility among ratifiable options and admitting mixed acts so that such
options exist; \textcite{richter1984}, \textcite{weirich1985}, and \textcite{rabinowicz1988} studied the
instability that makes every option look worse once chosen; and \textcite{skyrms1990dynamics}
modelled deliberation as a dynamical process whose rest points are deliberational equilibria,
connecting ratifiability with equilibrium in games \autocite{skyrms1990ratif}. \textcite{sobel1994}
collected much of this work.

\textcite{egan2007} reopened the debate with cases in which an act is evidence about the very causal
structure that determines its outcome---shooting as evidence of a lesion that would make one miss,
pressing a button that kills all psychopaths as evidence that one is a psychopath---and in which
causal decision theory seems to give the wrong verdict. The responses define the current landscape.
Deliberational theories recommend equilibria of a process of self-updating
\autocite{arntzenius2008,joyce2012,joyce2018}; benchmark theory assesses each option against a
state-dependent benchmark \autocite{wedgwood2013}; \textcite{spohn2012} argued that a properly
formulated causal theory one-boxes; \textcite{ahmed2012,ahmed2014book,ahmed2014dicing} revived the
evidential theory with cases in which causal reasoning fails; \textcite{harehedden2016} classified
decisions as self-reinforcing or self-frustrating and criticised the stability-based remedies;
\textcite{gallow2020} proposed a causal theory that manages the news an act carries; and
\textcite{gallow2026} argues for comparativism, on which instrumental value is fundamentally a matter
of pairwise comparison. The argument from wealth was reconsidered as well: \textcite{bales2018} and
\textcite{spencerwells2019} argued that it compares agents who face different problems, and
\textcite{wells2019} constructed a problem with equal opportunities in which causal reasoners
predictably end up richer. The collection edited by \textcite{ahmed2018ed} surveys the state of play.

\subsection{Policies, plans, and fairness}

A third line of work evaluates not acts but policies. Kuhn's analysis of extensive games introduced
perfect recall and the equivalence of mixed and behavioural strategies under it
\autocite{kuhn1953}. \textcite{hammond1988} derived expected utility from consequentialism and dynamic
consistency in decision trees; \textcite{mcclennen1990} defended resolute choice, the execution of an
ex-ante plan even when it is not optimal ex post; and \textcite{gauthier1986,gauthier1994} argued that
dispositions which are rational to adopt can make it rational to carry out assurances and threats.
\textcite{piccionerubinstein1997} separated planning-stage from action-stage optimality in the
absent-minded driver, and \textcite{aumannetal1997} showed that the planning-optimal strategy is
action-optimal although the converse fails; \textcite{arntzeniusetal2004} described sequences of
individually rational choices whose sum is disastrous and from which only binding escapes.

In artificial intelligence the policy view became a family of decision theories.
Updateless decision theory \autocite{dai2009} selects the input--output map that is optimal under
the agent's prior; timeless decision theory \autocite{yudkowsky2010} treats a decision as the output
of an abstract computation; and \textcite{soaresfallenstein2015} and \textcite{yudkowskysoares2017}
developed these ideas into functional decision theory, with a peer-reviewed presentation by
\textcite{levinsteinsoares2020}. Functional decision theory is claimed to do well on \emph{fair}
problems, those in which the outcome depends only on the agent's behaviour in the problem at hand;
its authors leave a formal criterion of fairness as future work and suggest that a general-purpose
theory may match or outperform all rivals on fair problems if one can be found.
\textcite{macaskill2019} objected with the Bomb, a problem in which the policy-optimal action is to take
a box known to contain a bomb, violating the principle that under certainty one should choose the
best outcome. \textcite{oesterheldconitzer2021} showed how to extract money from causal decision
theorists by predictive offers; \textcite{everittetal2015} formulated sequential, act-based and
policy-based versions of the causal and evidential theories; \textcite{belletal2021} showed that
value-based reinforcement learners in Newcomb-like environments converge, if at all, to ratifiable
policies, and that continuous environments always admit one; \textcite{halpern2020} gave conditions for
the existence of ratifiable mixed acts; \textcite{macdermottetal2023} characterised the theories with
mechanised causal graphs; and \textcite{oesterheldconitzer2024dese} and \textcite{cooperetal2024}
studied when self-locating beliefs make act-level causal or evidential reasoning agree with the
ex-ante optimal policy. \textcite{wolpertbenford2013} argued that the two Newcomb arguments presuppose
incompatible formalisations of how choice and prediction are related.

\subsection{Programs, proofs, and logical counterfactuals}

The functional theories treat a decision as the output of a computation, and so inherit a question
that the earlier theories could leave implicit: what would have happened had that computation
produced a different output? An idea of this kind is already in Hofstadter's \emph{superrationality},
on which agents who know that they reason identically should cooperate in a symmetric prisoner's
dilemma \autocite{hofstadter1983}. Game theory made programs the objects of choice.
\textcite{tennenholtz2004} introduced program equilibrium, in which players submit programs that can
read one another's code, and \textcite{barasz2014} showed that agents defined in provability logic
cooperate robustly with one another, while \textcite{critch2019} proved a resource-bounded version of
Löb's theorem underlying such results.

Within the functional programme, \textcite{soaresfallenstein2015} identified reasoning about
counterpossibles---conditionals whose antecedent is a non-actual output of a deterministic
computation---as the central open problem, and gave a proof-based formalisation that queries a
provability oracle. Such agents can act on spurious counterfactuals, provable only because the agent's own action is
provably different. The difficulty was raised in discussion of timeless decision theory
\autocite{drescher2009}, given an explicit Löbian form for proof-searching agents
\autocite{slepnev2010,slepnev2012}, and analysed by \textcite{bensontilsen2014}; it is now known as the
\enquote{five-and-ten} problem \autocite{fallenstein2013fiveandten,demskigarrabrant2019}. The
\enquote{chicken rule}, on which an agent takes any action it can prove it will not take, was
proposed to block such proofs \autocite{slepnev2011}. Fallenstein developed
modal versions of updateless decision theory, proved that for every provably extensional problem
some sound extension of arithmetic makes it optimal \autocite{fallenstein2014optimality,
fallenstein2014improving}, constructed problems on which every fixed modal theory fails
\autocite{fallenstein2014evil}, and proposed a partial semantics of \enquote{third-person
counterfactuals} \autocite{fallenstein2015third,fallenstein2015chicken}. Logical induction
\autocite{garrabrant2016li} gave a computable criterion for assigning probabilities to logical
sentences; its use for decisions requires exploration and meets obstacles of its own
\autocite{garrabrant2017obstacles}. \textcite{demskigarrabrant2019} placed these problems within a
wider account of agents embedded in the world they reason about, and \textcite{taylor2020}
proposed alternatives to logical counterfactuals altogether. In general reinforcement learning,
\textcite{leikehutter2015} showed that natural notions of optimality for universal agents are either
dependent on the prior or trivial. Two classical results of computability theory frame all of this:
Rice's theorem \autocite{rice1953}, that every non-trivial property of the functions that programs
compute is undecidable, and Kleene's second recursion theorem \autocite{kleene1938,moschovakis2010}, which
allows a program to refer to its own code.

In philosophy, the same conditionals are \emph{counterpossibles}. On the orthodox semantics of
\textcite{lewis1973counterfactuals} they are all vacuously true; non-vacuists evaluate them at impossible
worlds \autocite{nolan1997,brogaardsalerno2013}; and the exchange between \textcite{williamson2015,williamson2018}
and \textcite{bertoetal2018} concerns whether a semantics can treat them non-vacuously while respecting
the identifications that ordinary counterfactuals respect. \textcite{lewis1976} proved the classical
triviality theorem for a related identification, of conditional probabilities with probabilities of
conditionals.

\subsection{Impossibility results}

Four results are of the form this paper pursues. \textcite{skalse2021} proposed a general
counterexample to any decision theory: an opaque box that is empty exactly when the target theory
permits taking it. \textcite{fallenstein2014evil} proved the corresponding result in provability logic:
every fixed modal decision theory fails on some provably extensional problem. \textcite{briggs2010} proved that no completely general decision theory satisfies
both a Pareto condition across the evaluations of an agent's possible post-decision selves and a
self-sovereignty condition on which the comparison of two options depends only on how those two
options evaluate all options, modelling decision-theoretic paradoxes on voting paradoxes and Sen's
liberal paradox. \textcite{podgorski2022} proved that, for theories assigning each option an overall
value, independence of irrelevant alternatives forces the comparison of two options to depend only
on their values by their own lights, and so singles out the evidential theory; his tournament decision
theory escapes the result by producing no overall value.

These results draw on the theory of social choice, whose methods the present paper also uses.
\textcite{arrow1951} established that no aggregation rule satisfies a handful of individually
innocuous conditions on an unrestricted domain. The contraction condition on choice functions goes
back to \textcite{chernoff1954}; \textcite{arrow1959} showed that choice from all finite subsets is
rationalised by an ordering exactly when it satisfies the weak axiom of revealed preference; and
\textcite{sen1970,sen1971,sen1977,sen1993} named and analysed the conditions $\alpha$, $\beta$, and $\gamma$
and the limits of internal consistency of choice. Contraction consistency conflicts with majority
choice through the binary cycle, a fact central to the theory of tournament solutions
\autocite{moulin1986,laslier1997,brandtharrenstein2011,brandtbrillharrenstein2016}, and it continues to
generate impossibilities \autocite{blairetal1976,basumallikbhattacharya2026}.

\subsection{The question and the answer}

The thesis this paper set out to prove was that no decision theory is optimal on all fair decision
problems, the obstruction coming from mutually incompatible criteria of optimality rather than from
any demand for a general rule. Whether it is true depends on what a decision theory is required to
be, and the paper answers the question three times, for three successively stronger requirements.
Each strengthening moves the obstruction, and at each stage the obstruction turns out to lie in the
specification of the agent's alternatives rather than in the choice among them.

\emph{A theory as an assignment} (Part~\ref{part:single}). Section~\ref{sec:model} defines a class of
finite problems with prediction in which every counterfactual the criteria need is supplied by a
structural representation: for every policy, the law of the world's response is part of the problem.
A decision theory is any assignment of policies to problems, lookup tables included. The thesis is
then false: choosing an ex-ante maximiser in each problem is a decision theory that is ex-ante optimal
everywhere (Theorem~\ref{thm:universal}). The answer is cheap because the model hands the agent exactly
what causal, evidential, and functional theories disagree about. What remains true is shallow. Among
ex-ante optimality, structural causal dominance, ratifiability, commitment stability, and an evidential
criterion, every inconsistency has a core of one or two criteria (Theorems~\ref{thm:map15}
and~\ref{thm:map16}), the central pair being Nozick's. A conflict between ex-ante optimality and
dominance is possible exactly when the state law depends on the prescribed action
(Theorem~\ref{thm:characterisation}); it survives imperfect prediction and randomisation, with an exact
trade-off between the two kinds of regret (Proposition~\ref{prop:tradeoff}). Every within-problem
conflict reduces to the choice of the law under which an alternative is evaluated: the law induced by
the alternative, or the law induced by the actual choice.

\emph{A theory as a rule} (Part~\ref{part:across}). Requiring one rule across problems needs a common
description of them. Each problem is reduced to its causal matrix, whose rows are the belief states
induced by the agent's possible policies and whose diagonal is ex-ante value; every real square matrix
arises (Theorem~\ref{thm:realise}). With common affine invariance, the cross-problem axioms again
reduce to singleton and pairwise cores, one of which is Briggs's theorem (Theorem~\ref{thm:pairwise}).
With invariance under row translations---the matrix form of evaluating options against the baseline
of each belief state---a genuine four-axiom impossibility appears: neutrality, row-translation
invariance, and contraction force a rule to choose every option (Theorem~\ref{thm:circulant}), on
exactly those domains satisfying a reachability condition (Theorem~\ref{thm:DR}). Ex-ante evaluation,
in turn, is characterised by principal binary independence together with neutrality, affine
invariance, ordering, non-triviality, and monotonicity (Theorem~\ref{thm:selfesteem} and
Corollary~\ref{cor:three}), whose core is Podgorski's theorem. Both results depend on how a smaller
menu is modelled. If removing options also removes the beliefs that those options would induce, the
results hold; if the belief model is held fixed, row-baseline-invariant evaluation becomes menu
coherent, and contraction no longer forces ex-ante evaluation (Section~\ref{sec:menus} and
Theorem~\ref{thm:lambda}). The impossibilities concern the preservation of comparisons when the belief
perspectives associated with removed options disappear, not the preservation of choices when options
disappear.

\emph{A theory as a procedure} (Part~\ref{part:programs}). Finally, the supplied dependence is removed.
An agent is a computable function that receives the source code of its environment and must itself
determine what would happen under each of its outputs. With a designated substitution interface and a
termination promise, one algorithm is optimal everywhere (Theorem~\ref{thm:promiseopt}), the uniform
counterpart of Theorem~\ref{thm:universal}. The negative results separate into three kinds that have
nothing to do with counterfactual inference: queries about other problems
(Proposition~\ref{prop:foreign}); fixed copies of the agent, which defeat every agent that halts on all
inputs (Theorem~\ref{thm:fixedcopy}, Skalse's construction made well defined); and the uncomputability
of exact selection among computable expected values, which holds with no predictor at all
(Theorem~\ref{thm:exactsel}). The genuinely counterfactual question is how an agent should identify the
parts of a closed program that depend on its output. Literal, local, and global matching of the agent's
code each fail in a characteristic way (Theorems~\ref{thm:localmatch}--\ref{thm:haltred}), and
proof-based counterfactuals are spurious or undefined for ordinary programs (Theorem~\ref{thm:fiveten}
and Proposition~\ref{prop:guardedarith}). Axiomatising counterfactual operators gives a complete
consistency map (Theorem~\ref{thm:cfmap}). Calibration on independent problems, together with an
extensionality that identifies computations by their actual behaviour, forces causal evaluation
(Lemma~\ref{lem:projection}), and with ratifiability it forms a minimal inconsistent triple whose
witness is Death in Damascus (Theorem~\ref{thm:calibration}). Under an extensionality that compares
computations under every substitution, a single operator, ex-ante evaluation, satisfies every
desideratum (Corollary~\ref{cor:parametric}). The two evaluations of Part~\ref{part:single} thus
reappear as two notions of when computations are the same.

The conclusion of the paper unifies the three answers. \emph{No impossibility found here is an
impossibility of choosing well among specified alternatives. Each concerns how the alternatives are
specified: under which law an alternative is evaluated, which belief perspectives survive when options
are removed, and how a change in the agent's output is transported through the computations that
depend on it.}

Several results are known and are presented as such: the Newcomb conflict (Nozick), the absence of
ratifiable pure acts (Gibbard and Harper; Jeffrey), mixed existence under continuity (Halpern; Bell
and co-authors), the Pareto--self-sovereignty conflict (Briggs), the core of the binary
characterisation (Podgorski), the fixed-copy diagonal (Skalse; Fallenstein in modal form), the
spurious-proof mechanism behind five-and-ten, modal optimality (Fallenstein), and the recursion-theoretic
ingredients of the exact-selection theorem. The consistency maps, the characterisation of when
dominance and ex-ante optimality can conflict, the trade-off identity, the analysis of noisy
commitment, the circulant theorem and its exact domain criterion, the comparison of menu semantics,
the cardinality results on families of rules, the matching theorems, and the consistency map for
counterfactual operators were not found in the literature searched; their ingredients are classical,
and no claim of priority is made.

Part~\ref{part:single} (Sections~\ref{sec:model}--\ref{sec:robust}) treats single problems: the model,
the criteria and the possibility theorem, the witness problems, commitment and the within-problem
maps, and robustness. Part~\ref{part:across} (Sections~\ref{sec:matrix}--\ref{sec:mixed}) treats rules:
the causal matrix, the pairwise reduction, row translations and contraction, menu semantics, the
characterisation of ex-ante evaluation, problem types, and mixed policies. Part~\ref{part:programs}
(Sections~\ref{sec:programs}--\ref{sec:counterpossibles}) treats procedures: problems as programs,
computable optimisation and its limits, the inference of dependence from source, the axioms for
counterfactual operators, and the witness problems revisited together with the literature on
counterpossibles. Section~\ref{sec:conclusion} states the final thesis and the questions left open, and
Appendix~\ref{app:stipulations} lists every stipulation made.
